\documentclass[10pt,a4paper]{amsart}
\usepackage[margin=19mm]{geometry}
\usepackage{amsmath,amssymb,mathtools}
\usepackage[scr=rsfs,cal=euler]{mathalfa}
\usepackage{xfrac}
\usepackage[unicode,hidelinks,bookmarks=false]{hyperref}
\newcommand{\ie}{i.e.\ }
\newcommand{\cf}{cf.\ }
\newcommand{\resp}{respectively\ }
\newcommand{\Subsection}{\subsection}
\providecommand{\nobreakdash}{\nobreak\hbox{-}}
\setlength{\emergencystretch}{2em}
\multlinegap0pt
\usepackage{amssymb}
\usepackage{xcolor}
\usepackage{float}
\newtheorem{lemma}{Lemma}
\title{On the Bergman number of hyperbolic domains}
\author{Dimitrios Betsakos, Nikolaos Karamanlis}
\date{Source: arXiv:2607.14265v2 (CC BY 4.0)}
\begin{document}
\maketitle

\noindent\emph{Source and licence.} On the Bergman number of hyperbolic domains. Version arXiv:2607.14265v2 (CC BY 4.0), distributed by arXiv under the Creative Commons Attribution 4.0 International licence, http://creativecommons.org/licenses/by/4.0/, as recorded in the arXiv licence metadata for this exact version and shown on the abstract page https://arxiv.org/abs/2607.14265v2. Full licence notice: \texttt{licenses/arxiv-2607.14265-CC-BY-4.0.txt}.\par\smallskip

\noindent\textit{Editorial scope.} Counted unit: the article's lemma L1 (source0.tex lines 351--353). The article's complete original proof of it is delivered with it (source0.tex lines 354--408, one \texttt{proof} environment of 55 lines). No mathematics is written, repaired or re-derived: the statement and the proof are the article's own lines, in the article's own notation, and no retained line is substituted or re-flowed.  Retained uncounted as the source-local context the proof needs: lines 233--235.  Nothing was omitted from any retained span, so the package carries no omission record.  No retained line contains a citation, so the wrapper carries no bibliography and no bibliography entry.  No retained line contains a figure environment, an image-inclusion command or drawing code, so no image is delivered and the package declares no asset dependency.  Editorial formatting only, in addition: an AMS \texttt{amsart} preamble that restates the article's own declarations and macros used by the retained lines, namely the environments \texttt{lemma} on separate counters, together with the standard packages those lines need.  The retained environments print in this document as Lemma 1 in order of appearance, whereas the article prints the same items with the numbering of its own full document.  The complete native source of the article as distributed in the arXiv e-print for this exact version is bundled unchanged as \texttt{sources/205277/source0.tex}.
\begin{equation}\label{upperlambda}
	\lambda_D(z)\leq\frac{2}{\delta_D(z)}, \ z\in D.
\end{equation}

\begin{lemma}\label{L1}
The domain $G_\epsilon$ satisfies $b(G_\epsilon)<\infty$, for all $\epsilon>0$ sufficiently small.
\end{lemma}

\begin{proof}
We first observe that 
\[
b(G_\epsilon)=\liminf_{R\to\infty}\frac{\rho_{G_\epsilon}(0,C_R)}{\log R}\leq \liminf_{n\to \infty}\frac{\rho_{G_\epsilon}(0,C_{R_n})}{\log R_n}\leq \liminf_{n\to \infty}\frac{\ell_{G_\epsilon}([0,R_n])}{\log R_n}.
\]
By the Stolz-Ces\`{a}ro theorem, the last term in this inequality is no larger than
\[
\limsup_{n\to \infty}\frac{\ell_{G_\epsilon}([0,R_{n+1}])-\ell_{G_\epsilon}([0,R_n])}{\log R_{n+1}-\log R_n}=\limsup_{n\to \infty}\frac{\ell_{G_\epsilon}([R_n,R_{n+1}])}{\log R_{n+1}-\log R_n}.
\]
Hence
\begin{equation}\label{L1Eq1}
	b(G_\epsilon)\leq \limsup_{n\to \infty}\frac{\ell_{G_\epsilon}([R_n,R_{n+1}])}{\log R_{n+1}-\log R_n},
\end{equation}
which shows that it suffices to bound $\ell_{G_\epsilon}([R_n,R_{n+1}])$. Note that by \eqref{upperlambda},
\[
\ell_{G_\epsilon}([R_n,R_{n+1}]) \lesssim \int_{R_n}^{R_{n+1}}\frac{dt}{\delta_{G_\epsilon}(t)}.
\]
An easy calculation shows that there exists a number $M_n\in (R_n,R_{n+1})$ which satisfies $\frac{M_n}{R_{n+1}}\to 1/2$, as $n\to\infty$, and is equidistant from the points $R_n e^{i\pi R_n^{-2p}}$ and $R_{n+1}e^{i\pi R_{n+1}^{-2p}}$. For $\epsilon>0$ sufficiently small, we then have
\[\delta_{G_\epsilon}(t)= \begin{cases} 
	\big| t-R_n e^{i\pi R_n^{-2p}}\big|, & t\in (R_n,M_n) \\
	\big| t-R_{n+1} e^{i\pi R_{n+1}^{-2p}}\big|, & t\in (M_n,R_{n+1}).
\end{cases}
\]
Set $c_n=\cos\pi R_n^{-2p}$ and $s_n=\sin\pi R_n^{-2p}$. We split the integral $\int_{R_n}^{R_{n+1}}\frac{dt}{\delta_{G_\epsilon}(t)}$ into three pieces over the intervals  $(R_n,M_n)$, $(M_n,R_{n+1}c_{n+1})$,  and \\ $(R_{n+1}c_{n+1},R_{n+1})$.

Using the estimates $1-c_n\sim \frac{\pi^2}{2} R_n^{-4p}$ and $s_n \sim \pi R_n^{-2p}$, we find
\begin{align*}
	\int_{M_n}^{R_{n+1}c_{n+1}}\frac{dt}{\delta_{G_\epsilon}(t)} &\lesssim \int_{M_n}^{R_{n+1}c_{n+1}}\frac{dt}{R_{n+1}c_{n+1}-t+R_{n+1}s_{n+1}}\\
	&=\log\frac{R_{n+1}c_{n+1}-M_n+R_{n+1}s_{n+1}}{R_{n+1}s_{n+1}} \\
	&=\log\left(1+\frac{c_{n+1}-M_n/R_{n+1}}{s_{n+1}}\right) \\
	&\lesssim \log R_{n+1}.
\end{align*}
Furthermore, 
\begin{align*}
\int_{R_{n+1}c_{n+1}}^{R_{n+1}}\frac{dt}{\delta_{G_\epsilon}(t)}&\lesssim 
\int_{R_{n+1}c_{n+1}}^{R_{n+1}}\frac{dt}{t-R_{n+1}c_{n+1}+R_{n+1}s_{n+1}}\\
&=\log\left(\frac{1-c_{n+1}}{s_{n+1}}+1\right)\to 0,
\end{align*}
as $n\to\infty$. In a similar fashion, we find
\begin{align*}
\int_{R_n}^{M_n}\frac{dt}{\delta_{G_\epsilon}(t)}\lesssim \int_{R_n}^{M_n}\frac{dt}{t-R_nc_n+R_ns_n}&=\log\frac{M_n-R_nc_n+R_ns_n}{R_n-R_nc_n+R_ns_n}\\
&\lesssim \log\frac{M_n-R_nc_n+R_n^{1-2p}}{R_n^{1-4p}+R_n^{1-2p}}\\
&=\log \frac{R_n^{2p}\left(M_n/R_n-c_n+R_n^{-2p}\right)}{R_n^{-2p}+1}\\
&\lesssim \log \left(R_n^{2p-1}M_n\right).
\end{align*}
Using the fact that $M_n<R_{n+1}$, we infer
\[
\int_{R_n}^{M_n}\frac{dt}{\delta_{G_\epsilon}(t)}\lesssim 2p\log R_n+\log\frac{R_{n+1}}{R_n}.
\]
Upon adding the estimates for the three integrals above, dividing by $\log R_{n+1}-\log R_n$, and using the fact that the latter quantity is no less than $\log R_n$ (since $R_{n+1}>R_n^2$), we find that 
\[
\limsup_{n\to\infty}\frac{\int_{R_n}^{R_{n+1}}\frac{dt}{\delta_{G_\epsilon}(t)}}{\log R_{n+1}-\log R_n}\leq 3+2p<\infty.
\]                            
The conclusion follows by \eqref{L1Eq1}.
\end{proof}

\end{document}
